\documentclass[12pt]{article}
\usepackage[T1]{fontenc}
\usepackage[utf8]{inputenc}
\usepackage{lmodern}
\usepackage{textcomp}
\usepackage{enumitem}
\usepackage{geometry}
\usepackage{graphicx}
\usepackage{amsmath, amssymb}
\geometry{margin=1in}
\begin{document}
\begin{center}
Philosophy - Undergraduate Level Assessment 4 \
Total Marks: 40 \
Answer all questions.
\end{center}

\section*{Multiple Choice (10 marks)}
\begin{enumerate}[label=\arabic*.]
\item Augustine identifies evil with:
\begin{enumerate}[label=(\alph*)]
    \item a test from God.
    \item God's punishment for sin.
    \item human nature.
    \item the influence of the devil.
    \item a unique force, opposed to goodness.
\end{enumerate}
\item All things that are spoiled are inedible. Timothy is spoiled. So, Timothy is inedible.
\begin{enumerate}[label=(\alph*)]
    \item Equivocation
    \item Attacking the Person (ad hominem)
    \item Fallacy of Division
    \item Questionable Cause
\end{enumerate}
\item Anscombe claims that an adequate moral psychology would include:
\begin{enumerate}[label=(\alph*)]
    \item a detailed understanding of societal norms and expectations.
    \item a comprehensive review of historical philosophical theories.
    \item the influence of personal experiences on moral decisions.
    \item analyses of concepts such as "action" and "intention."
    \item the integration of religious beliefs into moral decisions.
\end{enumerate}
\item Anscombe claims that the notion of moral obligation is derived from the concept of:
\begin{enumerate}[label=(\alph*)]
    \item preference.
    \item natural law.
    \item self-interest.
    \item maximizing utility.
    \item ethical relativism.
\end{enumerate}
\item What can murtis be translated as?
\begin{enumerate}[label=(\alph*)]
    \item Sacrifices
    \item Blessings
    \item Offerings
    \item Prayers
    \item Apparitions
\end{enumerate}
\end{enumerate}

\section*{True / False (10 marks)}
\textit{Answer True or False}
\begin{enumerate}[label=\arabic*.]
\setcounter{enumi}{5}
\item True or False: Defenders of meat-eating argue that humans are rational beings, which grants them greater moral consideration than non-rational animals.
\item True or False: In Sikhism, strict adherence to the Five Ks is considered the essential path to achieving liberation and spiritual freedom.
\item True or False: Dershowitz believes that rule utilitarianism is the best framework for justifying terrorism under any circumstances.
\item True or False: Zhuangzi describes the state of ziran as embodying compassion, reflecting a deep connection with others and the natural world.
\item True or False: The fallacy of amphiboly is often referred to as reification, which involves treating abstract concepts as concrete entities.
\end{enumerate}

\section*{Long Form Response (20 marks)}
\textit{Answer each question with a clear and complete explanation.}
\begin{enumerate}[label=\arabic*.]
\setcounter{enumi}{10}
\item Construct a complete truth table for the argument \textasciitilde{}S ∨ T, \textasciitilde{}S · U, \textasciitilde{}T ∨ U, and analyze its
validity. If invalid, provide a counterexample and discuss its implications.
\item Construct a complete truth table for the argument (G ≡ H) · \textasciitilde{}I, \textasciitilde{}G ∨ (\textasciitilde{}H ∨ I) / G. Analyze
its validity and provide a counterexample if it is invalid, explaining your reasoning.
\end{enumerate}

\vfill
\noindent\textit{End of Paper}
\end{document}